\subsection{Approximate units in stellar algebras}

DEFINITION 7. --- Let A be a normed algebra. An approximate unit of A is a filter base F on the unit ball of A such that, for every x in A, the filter bases xF and Fx on A converge to x, in other words:

\[
\lim _ {f, \mathfrak {F}} f x = \lim _ {f, \mathfrak {F}} x f = x.
\]

If A is a stellar algebra, an approximate unit F is said to be increasing if F is a filter base on \(A_{+}\) .

Let A be a stellar algebra. We denote by \(A_{+}^{\leqslant1}\) (resp. \(A_{+}^{<1}\) ) the set of positive elements of A of norm \(\leqslant 1\) (resp. of norm \(< 1\)); these are the Hermitian elements of A whose spectrum is contained in \([0,1]\) (respectively in \([0,1)\)).

PROPOSITION 17. --- Let A be a stellar algebra, let \(\mathfrak{F}\) be a filter base on \(\mathrm{A}_{+}^{\leqslant 1}\) . For \(\mathfrak{F}\) to be an increasing approximate unit of A, it is necessary and sufficient that one have

\[
\lim _ {f, \mathfrak {F}} f x = x
\]

for every positive element x of A.

The condition is obviously necessary; let us prove that it is sufficient. Let \(\widetilde{\mathrm{A}}\) be the stellar algebra deduced from A by adjoining a unit element. Let \(x\) be an element of A and let \(f \in \mathrm{A}_{+}^{\leqslant 1}\) . We have

\[
\begin{array}{c} \| f x - x \| ^ {2} = \| (f x - x) (f x - x) ^ {*} \| = \| (f - 1) x x ^ {*} (f - 1) \| \\ \leqslant \| (f - 1) x x ^ {*} \| \end{array}
\]

because \(\| f - 1\| \leqslant 1\) (Lemma 12, \(c\) ) of I, p.~116). We therefore have

\[
\limsup _ {f, \mathfrak {F}} \| f x - x \| ^ {2} \leqslant \limsup _ {f, \mathfrak {F}} \| f x x ^ {*} - x x ^ {*} \|.
\]

Since \(xx^{*}\) is positive (th. 2 of I, p.~118), the hypothesis implies that

\[
\limsup _ {f, \mathfrak {F}} \| f x - x \| ^ {2} \leqslant \| x x ^ {*} - x x ^ {*} \| = 0,
\]

so that

\[
\lim _ {f, \mathfrak {F}} f x = x.
\]

Since the involution of A is continuous and since \(\mathfrak{F}\) is a filter base on \(A_{h}\) , it follows that

\[
\lim _ {f, \mathfrak {F}} x f = \lim _ {f, \mathfrak {F}} (f ^ {*} x ^ {*}) ^ {*} = \lim _ {f, \mathfrak {F}} (f x ^ {*}) ^ {*} = (x ^ {*}) ^ {*} = x.
\]

The proposition follows from this.

PROPOSITION 18. --- Let A be a stellar algebra. The ordered set \(A_{+}^{<1}\) is directed to the right (E, III, p.~12, def. 7) and the filter of its sections (TG, I, p.~38, example 2) is an increasing approximate unit of A.

Let \(\widetilde{\mathsf{A}}\) be the stellar algebra deduced from A by adjoining a unit element denoted 1 (def. 4 of I, p.~106).

Let \(g\) be the function on \([0,1)\) into \(\mathbf{R}_+\) defined by \(g(t) = t(1 - t)^{-1}\) . This is a continuous increasing bijection, whose inverse bijection is given by \(t \mapsto 1 - (1 + t)^{-1}\) .

Let us prove that the ordered set \(A_{+}^{<1}\) is directed to the right. Let x and y be elements of \(A_{+}^{<1}\) . Since \(g(0)=0\) and since \(\mathrm{Sp}_{\mathrm{A}}^{\prime}(x)\) and \(\mathrm{Sp}_{\mathrm{A}}^{\prime}(y)\) are contained in \([0,1)\) , the elements \(g(x)\) and \(g(y)\) are defined; they are positive, so that \(g(x)+g(y)\geqslant0\) . We can therefore form the element \(z=g^{-1}(g(x)+g(y))\) of A. We have \(\mathrm{Sp}_{\mathrm{A}}^{\prime}(z)\subset[0,1)\) , and hence \(z\in A_{+}^{<1}\) .

We have \(0 \leqslant g(x) \leqslant g(z)\) , whence \(1 \leqslant 1 + g(x) \leqslant 1 + g(z)\) . Assertion \(b\) ) of Lemma 14 of I, p.~119 implies that \(1 + g(x)\) and \(1 + g(z)\) are invertible and that \((1 + g(z))^{-1} \leqslant (1 + g(x))^{-1}\) . Consequently, \(z = 1 - (1 + g(z))^{-1} \geqslant 1 - (1 + g(x))^{-1} = x\) . Similarly, we have \(z \geqslant y\) . Therefore, \(z\) majorizes \(x\) and \(y\) in \(\mathrm{A}_{+}^{<1}\) . The ordered set \(\mathrm{A}_{+}^{<1}\) is thus directed to the right. Let \(\mathfrak{F}\) denote the filter of its sections

Let \(x\) be a positive element of A. For every integer \(n \geqslant 1\) , let \(e_n = g^{-1}(nx)\) ; we have \(e_n \in \mathrm{A}_+^{<1}\) . Let \(h_n\) be the continuous function on \(\mathbf{R}_+\) defined for all \(t \in \mathbf{R}_+\) by \(h_n(t) = t^2(1 - g^{-1}(nt)) = t^2 / (1 + nt)\) . We have \(|h_n(t)| \leqslant t / n\) for all \(t \geqslant 0\) , and hence \(\| x(1 - e_n)x \| = \| h_n(x) \| \leqslant \| x \| / n\) . In particular, \(x(1 - e_n)x\) tends to 0 as \(n\) tends to infinity.

Let \(\varepsilon > 0\) be a real number. Let \(n\) be an integer such that \(\|x(1 - e_n)x\| < \varepsilon\) . For every \(f \in \mathrm{A}_+^{<1}\) such that \(f \geqslant e_n\) , we then have

\[
\begin{array}{r l} & {\| x - f x \| ^ {2} = \| (1 - f) x \| ^ {2} = \| ((1 - f) x) ^ {*} (1 - f) x \| = \| x ^ {*} (1 - f) ^ {2} x \|} \\ & {\qquad \qquad \qquad \qquad \qquad \qquad \qquad \qquad \qquad \qquad = \| x (1 - f) ^ {2} x \|.} \end{array}
\]

Additionally, since \(0 \leqslant f \leqslant 1\) , it follows that \((1-f)-(1-f)^{2} = (1-f)f \geqslant 0\) (Lemma 12, d) of I, p.~116), hence \((1-f)^{2} \leqslant 1-f\) . Since \(1-f \leqslant 1-e_{n}\) , it follows that \(0 \leqslant (1-f)^{2} \leqslant 1-e_{n}\) . By lemma 14, a) and c) of I, p.~119, it follows

\[
\left\| x - f x \right\| ^ {2} = \left\| x (1 - f) ^ {2} x \right\| \leqslant \left\| x (1 - e _ {n}) x \right\| <   \varepsilon .
\]

We therefore have \(\lim_{f,\mathfrak{F}}fx = x\) for all \(x \in A_{+}\) . The filter \(\mathfrak{F}\) is therefore an approximate unit of A by proposition 17.

